\bookchapter{Глава 26: Передовые технологии смартфонов сегодня}

\chapterintro{Закончив с «последними делами» старого телефона, вернёмся к современности. 2024–2026 годы — период концентрированного технологического взрыва в мобильной индустрии: AI из облака спустился в карман, спутниковая связь перешла от флагманов к среднему сегменту, складные экраны из «новинки» стали «повседневностью», а технологии батарей впервые за много лет совершили прорыв на уровне материалов. В этой главе обозовём эти разворачивающиеся на наших глазах передовые технологии, чтобы помочь вам решить, «пора ли менять телефон» и «на какой».}

\booksection{I. AI-смартфон: маркетинговый трюк или реальная революция?}

С 2024 года практически все производители Android на флагманах заявляют о «AI-телефонах», а Apple в iOS 18 представила Apple Intelligence. Но какой же телефон считается «AI-телефоном»? В отрасли сформировались примерно три критерия:

\begin{itemize}
  \item \textbf{Онлайн LLM (Large Language Model)}:локально на телефоне может работать языковая модель с \textasciitilde{}7B параметров (например, лёгкие версии Qwen, Llama, DeepSeek), способная выполнять суммирование текста, перевод, генерацию контента без зависимости от облака. Способ проверки: если AI-функции работают в режиме полёта — значит, это онлайновая модель.
  \item \textbf{Производительность NPU (TOPS)}:мощность NPU (нейронного процессорного блока) измеряется в TOPS (триллионов операций в секунду). Флагменские чипы 2024 года (Snapdragon 8 Gen3, Dimensity 9300) имеют NPU 40–70 TOPS; чипы нового поколения 2025–2026 годов преодолевают 100 TOPS — это аппаратный порог для запуска больших моделей на устройстве.
  \item \textbf{Сценарии AI Agent}:не только «вопрос-ответ», но и выполнение цепочек задач через несколько приложений — например, «собери фотографии с вчерашней встречи в альбом и поделись с Чжан Санем»; система автоматически вызывает Галерею, Файлы, WeChat и др. Это ключевой признак отличия AI-телефона от «телефона с AI-ассистентом».
\end{itemize}

\booksection{II. Сравнение AI-стратегий производителей: кто делает серьёзно, а кто — наклейка?}

AI-маршруты производителей заметно различаются; при выборе обращайте внимание на:

\begin{itemize}
  \item \textbf{Apple Intelligence (iOS 18+)}:акцент на интеграцию с системными приложениями — суммирование писем, приоритет уведомлений, стилистическая правка текста, AI-ретушь фото. Основная модель — онлайновая; сложные задачи отправляются в Private Cloud Compute. Преимущества: защита приватности и плавность системы; недостаток: медленный запуск в китайскоязычном сегменте, часть функций поддерживается только с 2025 года.
  \item \textbf{Google Gemini Nano}:встроен в Pixel 8/9; поддерживает Circle to Search, Magic Compose, суммирование аудиозаписей. Gemini Nano — настоящая онлайновая малая модель, дружественная к открытой экосистеме; эталонная реализация AI-телефона на Android.
  \item \textbf{Huawei Pangu Large Model}:интегрирован в HarmonyOS Xiaoyi; поддерживает онлайновую генерацию текста, AI-ретушь, AI-удаление (убрать прохожих), суммирование звонков. После HarmonyOS NEXT 5.0 — дальнейшая системная интеграция; AI-возможности глубоко слиты с ОС.
  \item \textbf{Xiaomi MiLM}:встроен в HyperOS; поддерживает AI-портрет, AI-суммирование, AI-художник, суперразрешение. Xiaomi идёт по маршруту «онлайн + облако»: лёгкие модели на устройстве, тяжёлые задачи — в облаке.
  \item \textbf{OPPO (Andes Large Model) и vivo (Blue Heart Large Model)}:глубоко интегрированы в ColorOS 15 и OriginOS 5; ключевые возможности — AI-удаление, AI-суммирование документов, AI-ответы по изображениям, поиск фото на естественном языке («найди фото с прошлого года на море»).
\end{itemize}

\textbf{Рекомендация при выборе}:если вы часто используете AI для обработки рабочих документов и наведения порядка в информации, выбирайте аппарат с высокой онлайновой производительностью NPU и системной интеграцией AI; если AI-ретушь фото нужна лишь изредка — хватит основного среднего сегмента.

\booksection{III. Спутниковая связь: от флагманского «трюка» к стандарту среднего сегмента}

Спутниковая связь — одно из самых реальных технологий, спустившихся вниз за эти два года. В 2022 году Huawei Mate 60 и iPhone 14 были первыми, и это оставалось эксклюзивом флагманов; к 2025 году спутниковая связь уже доступна в сегменте 2000–3000 юаней.

\begin{itemize}
  \item \textbf{Короткие сообщения BeiDou в среднем сегменте}:функция коротких сообщений BeiDou (до 1000 китайских символов + координаты) уже доступна в Redmi, realme, iQOO и других моделях среднего сегмента; абонентской платы нет, использование бесплатное. Подходит для походов и автопутешествий.
  \item \textbf{Спутниковые звонки Tiantong}:Huawei Mate 60/70, Honor Magic6 и др. поддерживают звонки через спутник Tiantong-1; можно совершать/принимать спутниковые звонки и отправлять спутниковые сообщения. Требуется подключить спутниковый тариф у оператора (\textasciitilde{}10–200 юаней/мес.); подходит для пользователей, часто бывающих в зонах без покрытия.
  \item \textbf{Спутниковые сообщения iPhone (SOS)}:iPhone 14 и новее поддерживают спутниковый SOS: можно отправить текстовое сообщение экстренным службам и поделиться координатами. 2 года обслуживания бесплатно, в отдельных регионах доступно продление за плату.
  \item \textbf{Практическая рекомендация}:городским пользователям функция почти не нужна, но любителям активного отдыха и командировочным в отдалённые районы настоятельно рекомендуется модель со спутниковой связью — в критический момент это может спасти жизнь.
\end{itemize}

\booksection{IV. Зрелость экосистемы складных экранов: можно ли использовать как основной телефон?}

Складные экраны прошли путь от «технологической новинки» до «пригодности в качестве основного аппарата». Чтобы оценить их зрелость, смотрим на три аспекта: аппаратная долговечность, адаптация приложений, экосистема взаимодействия.

\begin{itemize}
  \item \textbf{Аппаратная долговечность}:ресурс шарнира — обычно 500 000+ циклов (при 100 открытиях/день хватает на 13 лет); каплевидный шарнир уменьшает складку на экране; защита IPX8 уже стандарт. Внутренний экран из UTG-стекла + защитная плёнка делает экран менее «хрупким».
  \item \textbf{Адаптация приложений}:основные приложения (WeChat, Taobao, Douyin, Bilibili, Office, WPS) уже адаптированы под большой экран; однако \textasciitilde{}30\% «длиннохвостых» приложений не адаптированы и отображаются с растяжением или чёрными полосами. Перед покупкой рекомендуется протестировать ваши ключевые приложения.
  \item \textbf{Взаимодействие в полускладке}:при складывании до 90°–120° некоторые приложения поддерживают режим «контент вверху / управление внизу», удобный для видеозвонков, селфи и просмотра видео. Селфи при полускладке — уникальный опыт складных экранов.
  \item \textbf{Разделение экрана и мультизадачность}:на большом экране можно одновременно открыть два приложения в разделённом экране или даже три в плавающих окнах. Логика разделения экрана в HyperOS, ColorOS, OriginOS уже достаточно зрелая, но разделение экрана в стиле iPadOS на iPhone пока не доступно.
  \item \textbf{Цена}:горизонтальные складни (книжного типа) опустились от 10 000+ до 5000–7000 юаней; вертикальные (раскладушки) — до 4000–5000 юаней. В 2025 году отдельные производители выпускают раскладушки за \textasciitilde{}3000 юаней — складные экраны по-настоящему входят в массовый рынок.
\end{itemize}

\textbf{Рекомендация при выборе}:если вы много читаете, работаете с документами и смотрите видео — горизонтальный складень стоит рассмотреть; если важны портативность и социальная атрибутика — вертикальная раскладушка подойдёт лучше. Рекомендуется хотя бы 30 минут протестировать в офлайн-магазине перед решением.

\booksection{V. 6G и Wi-Fi 7: скорость и границы связи следующего поколения}

6G ещё не по-настоящему развернут, но соответствующие технологии уже в пути:

\begin{itemize}
  \item \textbf{Прогресс 6G}:3GPP планирует завершить стандарт 6G к 2027 году, коммерческое развертывание — около 2030-го. Теоретическая пиковая скорость 6G — до 1 Tbps (в 50 раз выше 5G), задержка до 0.1 мс; поддерживается «интеграция связи и зондирования» — телефон сможет работать как низкоточный радар. Сейчас производители на этапе прототипной валидации; до потребительских продуктов ещё далеко.
  \item \textbf{Wi-Fi 7 (802.11be)}:с 2024 года — стандарт флагманов; теоретическая пиковая скорость — 46 Gbps (в 4 раза выше Wi-Fi 6); поддерживает ширину канала 320 МГц, 4K-QAM, MLO (многоканальная передача). На практике с Wi-Fi 7-роутером скорость на флагмане достигает 2–5 Gbps, что в 2–3 раза быстрее Wi-Fi 6E. Если вы приобрели Wi-Fi 7-телефон, не забудьте подобрать Wi-Fi 7-роутер, чтобы раскрыть все возможности.
  \item \textbf{Bluetooth 5.4 и LE Audio}:новые модели 2024–2026 годов обычно поддерживают Bluetooth 5.4 и LE Audio (низкоэнергетическое аудио); можно одновременно подключить две пары наушников к одному источнику, а Auracast позволяет нескольким людям в общественном месте синхронно слушать один и тот же аудиоисточник.
\end{itemize}

\booksection{VI. Кремний-углеродный анод: автономность наконец совершила прорыв на уровне материалов}

Автономность смартфонов многие годы критиковали за «зарядку раз в день»; с 2024 года распространение кремний-углеродных анодов реально улучшило ситуацию.

\begin{itemize}
  \item \textbf{Принцип технологии}:теоретическая ёмкость традиционного графитового анода \textasciitilde{}372 м А·ч/г; кремния — до 4200 м А·ч/г (\textasciitilde{}10×), но при зарядке/разрядке кремний расширяется на 300\%, что вызывает разрушение. Кремний-углеродный анод с нанокомпозитной структурой Si-C ограничивает расширение кремния до допустимого уровня, повышая плотность энергии батареи на 15–30\%.
  \item \textbf{Реальные результаты}:флагманы 2024–2025 годов с кремний-углеродным анодом (Xiaomi 14 Ultra, Honor Magic 6 Ultimate, vivo X100 Pro и др.) вмещают батарею 5000–6000 м А·ч в корпус толщиной 7–8 мм — на 800–1000 м А·ч больше предыдущего поколения; при интенсивном использовании хватает на полтора дня.
  \item \textbf{Сочетание с быстрой зарядкой $\geq$100 Вт}:кремний-углеродный анод хорошо совместим с быстрой зарядкой; зарядка 100 Вт / 120 Вт за 17 минут уже не новость. В 2025 году отдельные модели поддерживают 150–240 Вт: до 80\% за 10 минут.
  \item \textbf{Рекомендация при выборе}:если вы интенсивный пользователь, выбирайте модель после 2024 года с комбинацией «кремний-углеродный анод + зарядка $\geq$100 Вт» — тревога за автономность значительно снизится.
\end{itemize}

\booksection{VII. Обзор главы: одно предложение для каждой технологии}

\begin{itemize}
  \item AI-телефон, три критерия: онлайновая LLM, производительность NPU в TOPS, сценарии AI Agent через несколько приложений.
  \item Спутниковая связь: короткие сообщения BeiDou в среднем сегменте; звонки Tiantong — пока флагманы.
  \item Складные экраны: аппаратная зрелость, адаптация приложений \textasciitilde{}70\%, полускладок/разделение экрана — уникальный опыт.
  \item 6G ещё в пути, Wi-Fi 7 уже развернут, Bluetooth LE Audio начинает распространяться.
  \item Кремний-углеродный анод + зарядка $\geq$100 Вт: двойной прорыв автономности и зарядки.
\end{itemize}

\begin{tcbpractice}{Попробуйте: составьте «реестр технических требований» для следующего телефона}
Возьмите лист бумаги и по трём градациям «обязательно / плюс / не нужно» перечислите ваши требования к AI, спутниковой связи, складному экрану, Wi-Fi 7, быстрой зарядке и ёмкости батареи. Тогда в следующий раз в офлайн-магазине вы сможете быстро отличить реальные потребности от навязанных продавцом «киллер-фич». Например:

\begin{itemize}
  \item Обязательно: батарея $\geq$5000 м А·ч, зарядка $\geq$80 Вт, NFC, IP68.
  \item Плюс: онлайновое AI-суммирование, Wi-Fi 7, короткие сообщения BeiDou.
  \item Не нужно: складной экран (привык к моноблоку), спутниковые звонки (не езжу в глушь), экстремальная зарядка 240 Вт.
\end{itemize}

\end{tcbpractice}

\begin{tcbsummary}{Итоги главы}
\begin{itemize}
  \item AI-телефон: три критерия — онлайновая LLM, NPU TOPS, сценарии AI Agent
  \item AI-стратегии производителей: Apple — приватность и интеграция; Google — онлайновая модель и открытость; Huawei/Xiaomi/OPPO/vivo — у каждого свои сильные стороны
  \item Спутниковая связь: короткие сообщения BeiDou в среднем сегменте; звонки Tiantong — пока флагманский эксклюзив
  \item Складные экраны: аппаратная часть зрела, адаптация приложений \textasciitilde{}70\%, полускладок и разделение экрана — уникальный опыт
  \item Стандарт 6G — к 2027 году; Wi-Fi 7 уже распространён; Bluetooth LE Audio начинает развёртывание
  \item Кремний-углеродный анод делает 5000–6000 м А·ч стандартом флагманов; тревога за автономность значительно снижается
\end{itemize}

\end{tcbsummary}
